\documentclass[12pt]{letter}
\usepackage[UTF8]{ctex}
\usepackage[hidelinks]{hyperref}

\begin{document}

\begin{letter}{Information Fusion 编辑部}

\opening{尊敬的编辑：}

我们谨向贵刊投稿原创研究论文《异构大语言模型协作的反馈驱动自适应 DAG 执行：机制与适用边界》（Feedback-Driven Adaptive DAG Execution for Heterogeneous Large Language Model Collaboration: Mechanisms and Applicability Boundaries），恳请审阅。

\textbf{与期刊范围的匹配。}本文研究多模型异构大语言模型在复杂任务上的协同执行，属于贵刊覆盖的多源/多过程信息融合、多分类器与决策系统、自适应与自我改进架构以及融合系统中的计算资源优化方向。我们将异构模型能力视为待融合的信息源，通过任务有向无环图（DAG）在节点级进行能力融合，由反馈驱动的自适应执行层在质量—成本—时延约束下调度融合过程，并在财务、数学与代码三个任务域给出适用边界的机制刻画。

\textbf{主要贡献。}（1）节点级异构能力融合：将整题级模型选择转化为阶段级模型分配，并区分条件能力与传播能力；（2）反馈驱动的预算感知自适应执行：对失败节点及其受影响后继进行局部、选择性、可审计的调整，归因消融显示主要增益来自恢复动作本身而非调度逻辑；（3）适用边界的机制刻画：通过严格配对基准、精确 Oracle、恢复覆盖与三域正负对照及外部方法比较，揭示能力匹配、接口表达能力和反馈可诊断性三个条件共同决定多模型 DAG 协同的有效性。

\textbf{声明。}本文为原创工作，未同时投稿其他期刊；所有实验在冻结协议下一次完成，无利益冲突；所有作者已同意投稿。

恳请安排审阅。如需补充材料（冻结协议、任务清单、审计文件），我们可随时提供。

\closing{此致敬礼，\\ 作者代表}

\end{letter}

\end{document}
